\documentclass[12pt]{article}
\usepackage{amsmath, amssymb, amsthm}
\usepackage[margin=1in]{geometry}

\title{Problem 187: Semiprimes}
\author{Project Euler Solution}
\date{}

\begin{document}
\maketitle

\section{Problem Statement}

A composite number $n$ is a \textbf{semiprime} if $n = p \cdot q$ for primes $p$ and $q$ (not necessarily distinct). How many semiprimes are there below $10^8$?

\section{Mathematical Analysis}

\subsection{Enumeration by Smaller Prime Factor}

Write each semiprime as $n = p \cdot q$ with $p \le q$ both prime. Then:
\[
\text{Count} = \sum_{\substack{p \text{ prime} \\ p^2 < 10^8}} \left( \pi\!\left(\left\lfloor \frac{10^8 - 1}{p} \right\rfloor\right) - \pi(p-1) \right)
\]

where $\pi(x) = |\{q \text{ prime} : q \le x\}|$.

\subsection{Bounds on $p$}

Since $p \le q$ and $p \cdot q < 10^8$, we need $p^2 < 10^8$, giving $p < 10^4$. So $p$ ranges over all primes below $10000$.

\subsection{Proof of Correctness}

\begin{proof}
Each semiprime $n = p \cdot q$ with $p \le q$ has a unique representation with $p$ as the smallest prime factor. In our sum, when processing prime $p$, we count all primes $q$ in the range $[p, \lfloor(10^8-1)/p\rfloor]$. This counts $n = pq$ exactly once.
\end{proof}

\subsection{Sieve of Eratosthenes}

We sieve all primes up to $\lfloor(10^8-1)/2\rfloor = 49{,}999{,}999$, which is the largest possible value of $q$ (when $p = 2$). A cumulative prime count array allows $O(1)$ evaluation of $\pi(x)$.

\section{Complexity}

\begin{itemize}
    \item \textbf{Time}: $O(N \log \log N)$ for the sieve, where $N = 5 \times 10^7$.
    \item \textbf{Space}: $O(N)$ for the sieve and prefix sums.
\end{itemize}

\section{Answer}

\[
\boxed{17427258}
\]

\end{document}
